\documentclass[12pt, letterpaper]{article}

% --- Paquetes Esenciales ---
\usepackage[utf8]{inputenc}
\usepackage[spanish]{babel}
\usepackage{geometry}
\geometry{top=2.5cm, bottom=2.5cm, left=2.5cm, right=2.5cm}
\usepackage{booktabs}   % Para tablas con diseño profesional
\usepackage{amsmath}    % Para la fórmula matemática de tu Nivel de Riesgo
\usepackage{graphicx}   % Para cuando metas tus gráficas de Pareto y Evolución

% --- Información del Reporte ---
\title{Reporte de KPIs y Estrategia Comercial \\ \large Análisis del Último Trimestre}
\author{Emmanuel}
\date{Mayo de 2026}

\begin{document}

\maketitle
\newpage

\section{Introducción y Marco Metodológico}
El presente reporte tiene como objetivo evaluar el desempeño comercial, la salud de la cartera de clientes
y la eficiencia operativa del portafolio de productos durante el último trimestre \textit{Febrero - Abril}.

Para garantizar una correcta interpretación de los datos y asegurar una lectura unificada, se define a continuación
el glosario de las métricas avanzadas utilizadas en este diagnóstico.

\subsection{Glosario}
\begin{description}
    \item[Alcance.] Alcance sobre la cuota mensual del periodo.
    \item[Clientes Ganados] Clientes que no contaban con consumo en el periodo anterior, pero que sí lo tuvieron en el periodo actual.
    \item[Clientes Perdidos] Clientes que contaban con consumo en el periodo anterior, pero que no lo tuvieron en el periodo actual.
    \item[Venta neta.] Venta neta = \textit{Ingresos por ventas - Notas de crédito}.
    \item[Clientes con consumo activo.] Clientes que han registrado al menos una compra durante el periodo en cuestión.
    \item[Utilidad bruta.] Cálculo de utilidad = \textit{Venta neta - Costo de Venta} bruta del periodo.
    \item[Margen bruto.]  Calculo de margen bruto = \textit{Utilidad bruta / Venta neta}.
    \item[Indice de Gini.] Medida estadística entre 0 y 1 (expresada aquí en porcentaje) que mide la desigualdad en la distribución.
        Un Gini cercano al 100\% en clientes indica que la venta está peligrosamente concentrada en muy pocos compradores. En productos, indica dependencia de pocos SKUs.
    \item[Alcance de Cartera.] Porcentaje de clientes de la base total registrada que mantienen un consumo activo. Calculado como \textit{Clientes con consumo activo / Total de clientes registrados}.
    \item[Nivel de Riesgo.] Modelo matemático propio diseñado para ponderar la vulnerabilidad comercial de una ruta o territorio, calculándose mediante la fórmula:
    \begin{equation}
        \text{Nivel de Riesgo} = \alpha \cdot \text{Gini} + \beta \cdot (1 - \text{Alcance de Cartera})
    \end{equation}
    Donde $\alpha$ y $\beta$ son coeficientes de sensibilidad ajustables según los objetivos del negocio.
    \item[Diferencia Mediana vs. Media.] Desviación porcentual entre el promedio simple y el dato central de la muestra. Una brecha amplia indica la presencia de valores atípicos, clientes gigantes o compras extraordinarias, que sesgan la percepción del promedio común.
    \item[Regla de Pareto.] Principio que establece que aproximadamente el 80\% de los efectos provienen del 20\% de las causas. En este contexto, se refiere a la concentración de ventas en un pequeño porcentaje de clientes o productos.
\end{description}


\newpage

\section{Desempeño de Ventas y Cobertura}
Se analiza el comportamiento de los ingresos y el cumplimiento de metas por ruta.


\section{Dinámica y Conteo de Clientes activos}


\section{Análisis de Margen y Utilidad Bruta}



\section{Estructura de Riesgo Comercial}
% Aquí van tus índices de Gini, Pareto y el cálculo de la fórmula de riesgo

\end{document}